% Source: Wald GR, physical PDF pp. 454--457 (printed pp. 441--444).
% Coverage: Appendix C.3, Killing Vector Fields; Proposition C.3.1 and
%           equations (C.3.1)--(C.3.16).
% Figures/tables/footnotes: no figures/tables/footnotes.
% Status: source-reviewed and compile-clean.
% Uncertainties: none.

\subsection{Killing Vector Fields}\label{sec:C.3}

If \(\phi_t:M\to M\) is a one-parameter group of isometries,
\(\phi_t^*g_{ab}=g_{ab}\), the vector field \(\xi^a\) which generates \(\phi_t\) is
called a \emph{Killing vector field}. As already remarked below Eq.~\eqref{eq:C.2.2},
the necessary and sufficient condition for \(\phi_t\) to be a group of isometries is
\(\Lie_\xi g_{ab}=0\). Thus, according to Eq.~\eqref{eq:C.2.16}, the necessary and
sufficient condition that \(\xi^a\) be a Killing field is that it satisfy
\emph{Killing's equation}
\begin{equation}
  \nabla_a\xi_b+\nabla_b\xi_a=0,
  \tag{C.3.1}\label{eq:C.3.1}
\end{equation}
where \(\nabla_a\) is the derivative operator associated with \(g_{ab}\).

One of the most useful properties of Killing vector fields is given in the following
proposition.

\begin{proposition}\label{prop:C.3.1}
Let \(\xi^a\) be a Killing vector field and let \(\gamma\) be a geodesic with tangent
\(u^a\). Then \(\xi_au^a\) is constant along \(\gamma\).
\end{proposition}

\noindent\emph{Proof.} We have
\begin{equation}
\begin{aligned}
  u^b\nabla_b(\xi_au^a)
  &=u^bu^a\nabla_b\xi_a+\xi_au^b\nabla_bu^a\\
  &=0,
\end{aligned}
  \tag{C.3.2}\label{eq:C.3.2}
\end{equation}
since the first term vanishes by Killing's equation \eqref{eq:C.3.1} and the second
term vanishes by the geodesic equation. \(\square\)

Since in general relativity timelike geodesics represent the spacetime motions of
freely falling particles and null geodesics represent the paths of light rays,
Proposition~C.3.1 can be interpreted as saying that every one-parameter family of
symmetries gives rise to a conserved quantity for particles and light rays. This
conserved quantity enables one to determine the gravitational redshift in stationary
spacetimes and is extremely useful for integrating the geodesic equation when
symmetries are present (see Section~6.3).

Another useful formula relates the second derivative of a Killing field to the Riemann
tensor. By definition of the Riemann tensor, we have
\begin{equation}
  \nabla_a\nabla_b\xi_c-\nabla_b\nabla_a\xi_c
  =R_{abc}{}^d\xi_d.
  \tag{C.3.3}\label{eq:C.3.3}
\end{equation}
On the other hand, by Killing's equation, we can rewrite Eq.~\eqref{eq:C.3.3} as
\begin{equation}
  \nabla_a\nabla_b\xi_c+\nabla_b\nabla_c\xi_a
  =R_{abc}{}^d\xi_d.
  \tag{C.3.4}\label{eq:C.3.4}
\end{equation}
If we write down the same equation with cyclic permutations of the indices \((abc)\),
and then add the \((abc)\) equation to the \((bca)\) equation and subtract the
\((cab)\) equation, we obtain
\begin{equation}
\begin{aligned}
  2\nabla_b\nabla_c\xi_a
  &=(R_{abc}{}^d+R_{bca}{}^d-R_{cab}{}^d)\xi_d\\
  &=-2R_{cab}{}^d\xi_d,
\end{aligned}
  \tag{C.3.5}\label{eq:C.3.5}
\end{equation}
where the symmetry property (3.2.14) of the Riemann tensor was used in the last
equality. Thus, for any Killing field \(\xi^a\), we obtain the formula
\begin{equation}
  \nabla_a\nabla_b\xi_c=-R_{bca}{}^d\xi_d.
  \tag{C.3.6}\label{eq:C.3.6}
\end{equation}

An important consequence of Eq.~\eqref{eq:C.3.6} is that a Killing field, \(\xi^a\),
is completely determined by the values of \(\xi^a\) and
\(L_{ab}\equiv\nabla_a\xi_b\) at any point \(p\in M\); namely, if we are given
\((\xi^a,L_{ab})\) at \(p\), then \((\xi^a,L_{ab})\) at any other point \(q\) is
determined by integration of the system of ordinary differential equations
\begin{equation}
  v^a\nabla_a\xi_b=v^aL_{ab},
  \tag{C.3.7}\label{eq:C.3.7}
\end{equation}
\begin{equation}
  v^a\nabla_aL_{bc}=-R_{bca}{}^d\xi_dv^a,
  \tag{C.3.8}\label{eq:C.3.8}
\end{equation}
along any curve connecting \(p\) and \(q\), where \(v^a\) denotes the tangent to the
curve. Immediate corollaries of this result are (i) if a Killing field and its
derivative vanish at a point, then the Killing field vanishes everywhere, and (ii) on
a manifold of dimension \(n\), there can be at most
\(n+n(n-1)/2=n(n+1)/2\) linearly independent Killing fields [and, thus, at most an
\(n(n+1)/2\)-parameter group of isometries], since this is the dimension of the
space of initial data for \((\xi^a,L_{ab})\).

It is worth noting that if we contract Eq.~\eqref{eq:C.3.6} over \(a\) and \(b\), we
find
\begin{equation}
  \nabla^a\nabla_a\xi_c=-R_c{}^d\xi_d.
  \tag{C.3.9}\label{eq:C.3.9}
\end{equation}
Thus, in a vacuum spacetime, \(R_c{}^d=0\), \(\xi^a\) satisfies the source-free
Maxwell equation (4.3.15) for a vector potential in the Lorenz gauge. (There is a
sign difference in the Ricci tensor term between Eqs.~(4.3.15) and~\eqref{eq:C.3.9},
so Maxwell's equation is not satisfied when \(R_{ab}\ne0\).) The Lorenz gauge
condition \(\nabla_a\xi^a=0\) is also satisfied because of Killing's equation, and
thus all Killing fields in vacuum spacetimes give rise to solutions of Maxwell's
equation. Some solutions of physical interest can be obtained in this way
(Wald 1974b).

In the case of a hypersurface orthogonal Killing vector field, \(\chi^a\), we can
obtain a simple formula for \(\nabla_a\chi_b\). By Frobenius's theorem B.3.2, there
exists a vector field \(v^a\) such that
\begin{equation}
  \nabla_a\chi_b=\nabla_{[a}\chi_{b]}=\chi_{[a}v_{b]}.
  \tag{C.3.10}\label{eq:C.3.10}
\end{equation}
Assuming that \(\chi^a\) is not null, we may choose \(v^a\) to be orthogonal to
\(\chi^a\). Contracting Eq.~\eqref{eq:C.3.10} with \(\chi^b\), we obtain
\begin{equation}
  \frac{1}{2}\nabla_a(\chi^b\chi_b)
  =-\frac{1}{2}v_a\chi^b\chi_b.
  \tag{C.3.11}\label{eq:C.3.11}
\end{equation}
Hence, by solving Eq.~\eqref{eq:C.3.11} for \(v^a\) and substituting the result in
Eq.~\eqref{eq:C.3.10}, we find that an arbitrary hypersurface orthogonal Killing
field \(\chi^a\) with \(\chi^a\chi_a\ne0\) satisfies
\begin{equation}
  \nabla_a\chi_b=-\chi_{[a}\nabla_{b]}\ln\lvert\chi^c\chi_c\rvert.
  \tag{C.3.12}\label{eq:C.3.12}
\end{equation}

Finally, we mention two generalizations of the notion of Killing vector fields.
First, a \emph{conformal isometry}, \(\phi\), on a manifold, \(M\), with metric,
\(g_{ab}\), is defined to be a diffeomorphism \(\phi:M\to M\) for which there is a
function \(\Omega\) such that \(\phi^*g_{ab}=\Omega^2g_{ab}\). (The fact that
\(\phi\) is a diffeomorphism implies that \(\Omega\) is nonvanishing. The case
\(\Omega=1\), of course, corresponds to an ordinary isometry.) The infinitesimal
generator, \(\psi^a\), of a one-parameter group, \(\phi_t\), of conformal isometries
is called a \emph{conformal Killing vector field}. Clearly, the Lie derivative of
\(g_{ab}\) with respect to \(\psi^a\) must be proportional to \(g_{ab}\). Thus,
\(\psi^a\) satisfies
\begin{equation}
  \nabla_a\psi_b+\nabla_b\psi_a=\alpha g_{ab},
  \tag{C.3.13}\label{eq:C.3.13}
\end{equation}
where \(\nabla_a\) is the derivative operator associated with \(g_{ab}\). Taking the
trace of Eq.~\eqref{eq:C.3.13}, we evaluate the function \(\alpha\), thus obtaining
\begin{equation}
  \nabla_a\psi_b+\nabla_b\psi_a
  =\frac{2}{n}(\nabla^c\psi_c)g_{ab},
  \tag{C.3.14}\label{eq:C.3.14}
\end{equation}
where \(n=\dim M\). Equation~\eqref{eq:C.3.14} is known as the \emph{conformal
Killing equation}.

In Proposition~C.3.1, we proved that for any geodesic with tangent \(u^a\) and for
any Killing field \(\xi^a\), the inner product, \(\xi_au^a\), is constant along the
geodesic. The same calculation for a conformal Killing field yields
\begin{equation}
  u^b\nabla_b(\psi_au^a)=\frac{1}{n}(\nabla^c\psi_c)u^bu_b.
  \tag{C.3.15}\label{eq:C.3.15}
\end{equation}
Thus, in general, \(\psi_au^a\) is \emph{not} constant along a geodesic. However,
for a null geodesic we have \(u^bu_b=0\), so the right-hand side of
Eq.~\eqref{eq:C.3.15} vanishes. Thus, conformal Killing fields give rise to constants
of motion for null geodesics.

The second generalization we mention of a Killing vector is that of a Killing tensor.
A \emph{Killing tensor field} of order \(m\) on a manifold \(M\) with derivative
operator \(\nabla_a\) is defined to be a totally symmetric \(m\)-index tensor field,
\(K_{a_1\ldots a_m}=K_{(a_1\ldots a_m)}\), which satisfies the equation
\begin{equation}
  \nabla_{(b}K_{a_1\ldots a_m)}=0.
  \tag{C.3.16}\label{eq:C.3.16}
\end{equation}
Although Eq.~\eqref{eq:C.3.16} is a natural generalization of Killing's equation
\eqref{eq:C.3.1}, it should be noted that (aside from Killing vectors or Killing
tensors formed from products of Killing vectors) Killing tensor fields do not arise in
any natural way from groups of diffeomorphisms of \(M\). However, Killing tensors
share with Killing vectors the property of giving rise to constants of the motion: A
repetition of the proof of Proposition~C.3.1 shows that for any geodesic \(\gamma\)
with tangent \(u^a\), the quantity
\(K_{a_1\ldots a_m}u^{a_1}\cdots u^{a_m}\) is constant along \(\gamma\). The Kerr
metric (see Chapter~12) possesses a nontrivial Killing tensor \(K_{ab}\), and the
constant of motion to which it gives rise (together with the constants obtained from
the two Killing vectors) enables one to obtain all the geodesics explicitly.
